% ============ 4. Datos y método preliminar (D4) ============
\section{Datos disponibles y aproximación metodológica preliminar}

\subsection{Fuentes de datos (modalidad datos abiertos)}

Las fuentes son \textbf{datos abiertos oficiales} publicados por el Ministerio
de Ciencia, Tecnología e Innovación en el portal nacional
\texttt{datos.gov.co} (plataforma Socrata). Ya fueron \textbf{abiertas y
verificadas por el equipo} (conteo real de filas, columnas y valores), no
supuestas:

\begin{longtable}{p{2.6cm}p{4.4cm}p{4.2cm}p{2.8cm}}
\caption{Descripción de las fuentes de datos}
\label{tab:fuentes}\\
\toprule
\textbf{Atributo} & \textbf{Padrón de investigadores} & \textbf{Producción de grupos} & \textbf{Observaciones} \\
\midrule
\endfirsthead
\toprule
\textbf{Atributo} & \textbf{Padrón} & \textbf{Producción} & \textbf{Observaciones} \\
\midrule
\endhead
Entidad & MinCiencias & MinCiencias & datos.gov.co \\
Identificador & \texttt{bqtm-4y2h} & \texttt{33dq-ab5a} & Enlace en referencias \\
Periodo & 2013--2021 & 2013--2021 & 6 ventanas \\
Unidad de observación & Investigador $\times$ convocatoria & Producto $\times$ autor & — \\
Volumen declarado & 77\,237 filas; 30\,086 únicos & 3\,166\,629 filas & Verificado: el XLSX del clon tiene 50\,891 \\
Variables & 30 (categoría, área OCDE, género, geografía, diversidad, filiación) & 15 (producto, categoría, tipo, grupo, autor) & Diccionario en el catálogo \\
Acceso & API Socrata / descarga CSV & API Socrata (paginada) & Sin autenticación \\
Licencia & Datos abiertos del portal (uso libre con atribución) & Ídem & Verificar enlace oficial \\
Llave de cruce & \texttt{id\_persona\_pr} & \texttt{id\_persona\_pd} & $\leftrightarrow$ + \texttt{id\_convocatoria} \\
\bottomrule
\end{longtable}

\subsection{Variables ligadas a cada objetivo}

\begin{itemize}
  \item \textbf{OE1 (calidad):} \texttt{ID\_CONVOCATORIA}, \texttt{ANO\_CONVO},
        \texttt{ID\_PERSONA\_PR} (unicidad), \texttt{EDAD\_ANOS\_PR} (validez),
        codificación del archivo (consistencia).
  \item \textbf{OE2 (longitudinal):} \texttt{ID\_PERSONA\_PR},
        \texttt{ANO\_CONVO}, \texttt{ID\_CLAS\_PR}/\texttt{NME\_CLASIFICACION\_PR}
        (transiciones), más covariables área y región.
  \item \textbf{OE3 (equidad):} \texttt{NME\_DEPARTAMENTO\_RES\_PR},
        \texttt{NME\_GENERO\_PR}, \texttt{TXT\_GRUPO\_ETNICO},
        \texttt{TXT\_POBLACION\_DISCA}, \texttt{ID\_VICTIMA\_CONFLICTO}.
\end{itemize}

\subsection{Ruta metodológica preliminar}

El diseño sigue un flujo en tres etapas alineado con los objetivos:
\textbf{(1) auditoría de calidad} del registro (perfilado por las dimensiones
de Wang \& Strong, contratos de esquema y normalización reproducible);
\textbf{(2) construcción del panel longitudinal} y estimación de transiciones
(matrices de Markov con test de la propiedad markoviana, complementadas con
supervivencia Kaplan--Meier/Cox); y \textbf{(3) medición de equidad} (HHI,
Gini y Theil con intervalos de confianza bootstrap para lo territorial; razones
de representación con IC frente al censo DANE 2018 para género, etnia y
discapacidad). Todo se implementa en Python sobre un almacén DuckDB, con el
código, los datos y los resultados versionados y reproducibles.

\textbf{Alternativa considerada y descartada:} se evaluó complementar con
indicadores de producción individual (índice $h$ y variantes) a partir del
dataset de producción; se descartó en esta fase porque el cruce de 3,2
millones de filas requiere descargar ~1,2\,GB y su disponibilidad es
intermitente, y porque el objetivo de esta entrega es la auditoría del padrón;
el análisis de productividad queda como extensión (se declara en el alcance).

\textbf{Plan si el dato no llega o cambia:} ambas fuentes tienen descarga
directa por API Socrata y respaldo versionado; si el dataset de producción no
estuviera disponible a tiempo, los objetivos OE1--OE3 no dependen de él (se
audita el padrón y su artefacto consolidado). Si cambiara el esquema del
padrón, los contratos de esquema detectarán el cambio y se documentará la
re-ingesta.

\subsection{¿Por qué no está ya resuelto con estos datos?}

Aunque los datos son públicos desde 2021 y existen análisis descriptivos, no se
ha publicado —hasta donde alcanza esta revisión— una \textbf{auditoría
estadística reproducible} del padrón que combine las tres dimensiones (calidad
+ longitudinal + equidad) con los métodos de la ciencia de la ciencia y que
documente explícitamente los defectos del registro (discrepancia 77\,237 vs.\
50\,891, atípicos, mojibake). Nuestro aporte es ese: no repetir un análisis
descriptivo, sino \textbf{auditar el instrumento} y entregar un marco de
indicadores reutilizable.
